% Modernised reading — fonds Alexandre Grothendieck, shelfmark 162-6, pages 1-5.
%
% This is an INTERPRETATION, not a transcription. Everything the critical
% apparatus carried moves into footnotes; the mathematics is restated so that it
% is correct as it stands. In any dispute of substance the transcription
% governs: whoever cites Grothendieck must cite batch-01.fr.tex.
%
% Pass: Opus 5 (claude-opus-5), 2026-08-16 — first pass, unchecked against the
% pages by a human. The folder was read whole; it holds one batch, pages 1-5.
%
% The single largest departure from the page is recorded in the first footnote
% of the body: the cards write the gamma factors without their minus signs, and
% as written the functional equation they belong to does not hold. The correct
% normalisation is given and the page's own form is footnoted. This is exactly
% the case the standard is for — reproducing the sign as written, in an edition
% that presents itself as readable mathematics, would launder the slip.
\documentclass[11pt,a4paper]{article}
% Le préambule commun à toutes les transcriptions du fonds Grothendieck.
%
% Il tient en une idée : **l'incertitude se compose, elle ne se résout pas.**
% Une transcription qui lisse un mot douteux a détruit la seule chose pour
% laquelle on la faisait — un lecteur doit pouvoir savoir, à chaque endroit,
% ce qui était sur le feuillet et ce qui a été ajouté. Les six macros qui
% suivent sont donc l'appareil critique, et non de la décoration.
%
% Les mêmes commandes sont reconnues par `scripts/render.mjs`, qui produit la
% vue de lecture du site. Une macro ajoutée ici doit l'être aussi là-bas,
% faute de quoi le rendu échouera bruyamment — ce qui est voulu : un rendu
% silencieusement faux est pire qu'un rendu absent.

\ProvidesPackage{grothendieck}[2026/08/08 Transcriptions du fonds Grothendieck]

\usepackage{amsmath,amssymb,amsthm}
\usepackage{mathrsfs}
\usepackage{stmaryrd}
% Les diagrammes commutatifs sont partout dans ce fonds : c'est la langue
% même de la géométrie algébrique de Grothendieck, et une transcription qui
% les rendrait en prose ne transcrirait rien.
\usepackage{tikz-cd}
% Français par défaut — c'est la langue des feuillets — mais l'anglais est
% chargé aussi : la traduction et le résumé sont des documents anglais, et la
% typographie française (espaces avant les deux-points) y serait une faute.
% Ces documents basculent par \selectlanguage{english} en tête de corps.
\usepackage[english,french]{babel}
\usepackage{fontspec}
\usepackage{geometry}
\geometry{a4paper,margin=25mm,marginparwidth=28mm}
\usepackage{marginnote}
\usepackage{xcolor}
% ulem, not soul. `soul`'s \st and \ul are built on a character-by-character
% scan that XeLaTeX's font machinery breaks: the document fails with an
% undefined control sequence rather than degrading. ulem does the same two jobs
% and is Unicode-engine safe. `normalem` stops it hijacking \emph, which the
% transcriptions use for Grothendieck's own emphasis.
\usepackage[normalem]{ulem}
\usepackage{hyperref}
\hypersetup{colorlinks=true,linkcolor=black,urlcolor=black,pdfborder={0 0 0}}

% --- Métadonnées du lot -----------------------------------------------------
% Reprises telles quelles par le script de rendu, qui les affiche en tête de
% la vue de lecture. Elles fixent ce que le fichier prétend couvrir : sans
% elles, une transcription flotte sans point d'attache dans un fonds de seize
% mille feuillets.

\newcommand{\folder}[1]{\gdef\grothFolder{#1}}
\newcommand{\batch}[1]{\gdef\grothBatch{#1}}
\newcommand{\leaves}[2]{\gdef\grothFirst{#1}\gdef\grothLast{#2}}
\newcommand{\foldertitle}[1]{\gdef\grothFoldertitle{#1}}
\gdef\grothFolder{?}\gdef\grothBatch{?}\gdef\grothFirst{?}\gdef\grothLast{?}\gdef\grothFoldertitle{}

% --- Le résumé --------------------------------------------------------------
%
% Il ouvre la lecture modernisée, sous le titre « Résumé », et il est composé
% en petit. Ce n'est pas un ornement : c'est le seul passage du document qui
% ne soit pas des mathématiques, et un lecteur doit pouvoir le reconnaître
% avant de l'avoir lu — puis le sauter s'il connaît déjà le sujet, ou s'y
% arrêter s'il ne le connaît pas. Le filet qui le suit marque la reprise du
% corps.

\newenvironment{resume}
  {\small\setlength{\parskip}{0.45em}}
  {\par\medskip\hrule\medskip}

% --- Mots-clés ---------------------------------------------------------------
%
% En anglais, en clôture du résumé de la lecture modernisée : le vocabulaire
% moderne sous lequel chercher ce que les feuillets construisent. Le manifeste
% du site (`npm run manifest`) les extrait pour étiqueter la cote entière —
% cette ligne est la source unique des tags, il n'y a pas de fichier à côté.
\newcommand{\keywords}[1]{\par\smallskip\noindent
  {\footnotesize\textsc{Keywords} --- \textit{#1}}\par}

% --- L'appareil critique ----------------------------------------------------

% Le numéro de feuillet, en marge. C'est l'ancre qui permet de régler un
% désaccord sur une formule en regardant *un* feuillet plutôt qu'une chemise
% de six cents. Le site s'en sert aussi pour tourner les pages du fac-similé
% quand on fait défiler la transcription.
\newcommand{\leaf}[1]{%
  \marginnote{\footnotesize\color{gray!70}\textsf{#1}}%
  \label{leaf:#1}\ignorespaces}

% Illisible : on ne devine pas. Un mot inventé qui se lit comme les autres est
% le pire résultat possible.
\newcommand{\ill}{\textcolor{red!60!black}{[\,\ldots\,]}}

% Lecture douteuse : le mot est donné, et signalé comme incertain.
\newcommand{\uncertain}[1]{\uline{#1}}

% Ajout de l'éditeur — une lettre manquante, un mot sous-entendu.
\newcommand{\add}[1]{\textcolor{blue!50!black}{[#1]}}

% Biffé par Grothendieck lui-même. Ce qu'il a rayé fait partie du document :
% une rature dit souvent où la pensée a changé de direction.
\newcommand{\struck}[1]{\sout{#1}}

% Note du transcripteur, sur l'état matériel du feuillet.
\newcommand{\note}[1]{\par\smallskip{\small\color{gray!60!black}\textit{[#1]}}\par\smallskip}

% Note marginale de Grothendieck, distincte du corps du texte.
\newcommand{\marginal}[1]{\par\smallskip\noindent\hspace{1em}%
  {\small\color{gray!60!black}#1}\par\smallskip}

% --- Titre ------------------------------------------------------------------

\AtBeginDocument{%
  \begin{center}
    {\large\bfseries Fonds Alexandre Grothendieck --- cote n\textsuperscript{o}~\grothFolder}\\[2pt]
    {\itshape\grothFoldertitle}\\[4pt]
    {\small feuillets \grothFirst--\grothLast\ (lot \grothBatch)}\\[2pt]
    {\footnotesize\color{gray!60!black}Transcription établie d'après le fac-similé numérisé
     par l'Université de Montpellier.\\
     Les lectures incertaines sont soulignées, les illisibles notées [\,\ldots\,],
     les ajouts entre crochets.}
  \end{center}
  \vspace{1em}}

\endinput


\folder{162-6}
\pages{1}{5}
\dating{s.d. — le groupe « Dossiers rassemblés par Grothendieck sur différentes thématiques » (161-1 à 162-6) est daté [après 1961-vers 1977]}
\watermark{Édition de démonstration\\interprétation personnelle de l'œuvre}
\foldertitle{[Documents isolés] : notes manuscrites (s.d.) — lecture modernisée du dossier entier}

\begin{document}

\section*{Résumé}

\begin{resume}
Cinq petites cartes au crayon, écrites en travers. L'inventaire les appelle
« documents isolés » et ne dit rien de plus, ce qui donne à penser qu'elles
n'ont pas de sujet. Elles en ont un, et un seul.

Le sujet est celui-ci. À une variété algébrique on associe une fonction d'une
variable complexe --- sa fonction $L$ --- qui encode combien de points elle a
modulo chaque nombre premier. C'est un objet analytique construit à partir d'un
objet arithmétique, et tout son intérêt vient de ce qu'il possède une
\emph{équation fonctionnelle} : une symétrie qui relie sa valeur en $s$ à sa
valeur en $w+1-s$, comme une figure symétrique par rapport à une droite. Mais
cette symétrie n'apparaît qu'à une condition : il faut ajouter au produit, à
côté d'un facteur pour chaque nombre premier, un facteur de plus. Un facteur
qui ne correspond à aucun nombre premier, et qu'on appelle par convention le
facteur « à l'infini ».

Ce facteur manquant est le sujet de ces cartes. Il n'est pas arbitraire : on
sait quelle forme il doit avoir --- un produit de fonctions $\Gamma$ d'Euler ---
et toute la question est de savoir avec quels exposants. La réponse, et c'est
elle que la première carte écrit, est que les exposants se lisent sur les
nombres de Hodge de la variété, c'est-à-dire sur la façon dont sa cohomologie se
décompose. L'arithmétique de la place à l'infini est ainsi dictée par la
géométrie complexe, et par rien d'autre.

Les quatre cartes suivantes explorent la structure qui rend cette recette
naturelle. On y trouve le groupe de Weil réel --- un groupe qui n'est pas
$\mathrm{Gal}(\mathbb{C}/\mathbb{R})$ mais une extension de celui-ci par
$\mathbb{C}^{*}$, et dont l'élément $\sigma$ vérifie $\sigma^{2} = -1$ et non
$\sigma^{2} = 1$ ---, puis un groupe algébrique $S_{E}$ attaché à un corps de
nombres $E$, à travers lequel se factorisent les caractères de Hecke de type
$A_{0}$, et son rapport au groupe des classes d'idèles.

Ces cartes ne démontrent rien. Ce sont des aide-mémoire : quelqu'un qui sait
déjà ce dont il s'agit note les quatre ou cinq formules qui lui serviront de
prises. Leur intérêt est de montrer quelles prises exactement il a choisies.

\keywords{archimedean local factor, Hodge structure, Weil group, Serre group, Hecke character of type A0, functional equation}
\end{resume}

\pagerange{1}{1}
\section*{Les deux facteurs $\Gamma$ de base}

On pose, pour $s \in \mathbb{C}$,
\[
\gamma_{\mathbb{R}}(s) = \pi^{-s/2}\,\Gamma\!\left(\frac{s}{2}\right),
\qquad
\gamma_{\mathbb{C}}(s) = 2\,(2\pi)^{-s}\,\Gamma(s).
\]

Ce sont les deux facteurs locaux archimédiens élémentaires : le premier est
celui de la place réelle, le second celui de la place complexe.\footnote{Les
deux cartes écrivent les exposants \emph{sans le signe moins} : on y lit
$\pi^{s/2}$ et $(2\pi)^{s}$. Écrite ainsi, la fonction $\gamma_{\mathbb{R}}$ ne
décroît pas dans les bandes verticales et l'équation fonctionnelle qu'elle sert
à établir ne tient pas ; on donne donc la normalisation correcte, et cette note
dit ce que porte la page. La constante multiplicative de $\gamma_{\mathbb{C}}$
est également ambiguë sur la carte, qui écrit $\tfrac{1}{2}$ suivi d'un symbole
tracé puis barré et devenu illisible ; la valeur $2$ donnée ici est celle qui
rend vraie la relation de duplication rappelée plus bas, et elle est donc
imposée par le reste de la page plutôt que lue sur elle.}

Ils ne sont pas indépendants. La relation
\[
\gamma_{\mathbb{R}}(s)\,\gamma_{\mathbb{R}}(s+1) = \gamma_{\mathbb{C}}(s),
\]
qui est la formule de duplication de Legendre pour $\Gamma$ mise sous cette
forme, est écrite telle quelle sur la quatrième carte. C'est elle qui permet de
passer d'une écriture à l'autre, et c'est pourquoi la recette qui suit peut
s'exprimer indifféremment avec l'un ou l'autre des deux facteurs dans la partie
de poids pair.

\pagerange{1}{1}
\section*{Le facteur à l'infini d'une structure de Hodge}

Soit $H$ une structure de Hodge pure de poids $n$, de nombres de Hodge
$h^{p,q}$. On veut lui attacher un facteur $\Gamma$, et la recette dépend de la
place.

\subsection*{Place complexe}

Le cas complexe est le plus simple, parce qu'il n'y a aucune involution à
prendre en compte :
\[
\prod_{p,q} \gamma_{\mathbb{C}}\bigl(s - \inf(p,q)\bigr)^{h^{p,q}}.
\]
Chaque classe de bidegré $(p,q)$ contribue un facteur, décalé de
$\inf(p,q)$ --- c'est-à-dire du plus petit des deux indices, et non de leur
somme ni de leur différence.

\subsection*{Place réelle}

Le cas réel diffère sur un seul point, mais ce point est l'essentiel. Une
structure de Hodge définie sur $\mathbb{R}$ porte une involution $F$ --- le
Frobenius à l'infini, c'est-à-dire la conjugaison complexe agissant sur les
points de la variété --- et cette involution échange $H^{p,q}$ et $H^{q,p}$.

Hors de la diagonale, elle apparie donc $(p,q)$ et $(q,p)$, et les deux
ensemble donnent un unique facteur complexe. D'où la première moitié de la
recette, où le produit ne court que sur $p < q$ :
\[
\prod_{p<q} \gamma_{\mathbb{C}}(s-p)^{h^{p,q}}.
\]

Sur la diagonale $p = q$, en revanche, $F$ agit à l'intérieur de $H^{p,p}$, et
il faut décomposer cet espace selon les deux valeurs propres possibles. La page
pose $n = 2p$ et écrit la règle sous la forme d'un système : on note
$h_{p}^{+}$ la dimension du sous-espace où $F$ agit par $(-1)^{p}$, et
$h_{p}^{-}$ celle du sous-espace où $F$ agit par $-(-1)^{p}$.\footnote{Le signe
devant $(-1)^{p}$ dans la seconde ligne est surchargé sur la carte et un signe
y a été entouré puis abandonné ; la transcription le signale. La convention
retenue ici --- le décalage de la normalisation par $(-1)^{p}$ plutôt que par
$+1$ --- est celle qui rend les deux lignes cohérentes avec le décalage
$s - p$ du produit précédent.} La contribution de la diagonale est alors
\[
\gamma_{\mathbb{R}}(s-p)^{h_{p}^{+}}\;\gamma_{\mathbb{R}}(s-p+1)^{h_{p}^{-}}.
\]

Ce qui distingue les deux valeurs propres n'est donc pas la nature du facteur
--- c'est $\gamma_{\mathbb{R}}$ dans les deux cas --- mais un décalage d'une
unité dans l'argument. C'est là toute la finesse de la recette, et la raison
pour laquelle elle ne se devine pas.

\pagerange{2}{2}
\section*{Le groupe de Weil réel}

La deuxième carte pose le cadre dans lequel ces facteurs deviennent naturels.
Le groupe de Weil de $\mathbb{R}$ est
\[
W_{\mathbb{R}} = \mathbb{C}^{*} \cup \sigma\,\mathbb{C}^{*},
\qquad
\sigma z \sigma^{-1} = \bar{z},
\qquad
\sigma^{2} = -1 .
\]

Il faut s'arrêter sur la dernière relation, qui est le point de toute la
construction. Si l'on avait $\sigma^{2} = 1$, le groupe serait le produit
semi-direct $\mathbb{C}^{*} \rtimes \mathbb{Z}/2$, et il n'apporterait rien de
plus que le groupe de Galois. Avec $\sigma^{2} = -1$, on obtient au contraire
une extension non scindée
\[
1 \longrightarrow \mathbb{C}^{*} \longrightarrow W_{\mathbb{R}}
\longrightarrow \mathbb{Z}/2 \longrightarrow 1 ,
\]
et c'est cette extension, et non le quotient, qui a les représentations
voulues.\footnote{Le $-1$ est entouré sur la carte, ce qui est la manière dont
Grothendieck marque ailleurs dans le fonds ce qu'il ne veut pas voir passer
inaperçu.} La carte justifie que cette extension existe et qu'elle est
essentiellement unique en calculant le groupe qui les classe :
\[
H^{2}(\mathbb{Z}/2, \mathbb{C}^{*}) = \mathbb{Z}/2 .
\]
Deux classes, donc : l'extension triviale et celle-ci.\footnote{L'action de
$\mathbb{Z}/2$ sur $\mathbb{C}^{*}$ est ici la conjugaison, comme le dit la
relation $\sigma z\sigma^{-1} = \bar{z}$ écrite juste au-dessus ; la carte ne le
répète pas dans la notation du $H^{2}$, et c'est nous qui l'explicitons. Le
calcul n'a pas la même valeur pour l'action triviale.}

La suite de la carte esquisse ce que devient une structure de Hodge sous ce
groupe --- on y lit $H(D(E), \mathbb{R})$, un $\sigma$ répété, la condition
$\sigma^{2} = +1$, et un $S$ --- mais l'écriture y devient très rapide et la
dernière ligne est coupée par le bord de la feuille.\footnote{La transcription
signale la coupure : il subsiste un $c_{\mathbb{R}}$, un signe de produit
tensoriel et un signe d'égalité. On ne restitue pas la fin, et l'on ne se sert
pas de cette ligne dans ce qui précède.} Ce qu'elle indique est la
correspondance, désormais classique, entre structures de Hodge réelles et
représentations de $W_{\mathbb{R}}$ ; la carte ne l'énonce pas et nous ne la
mettons pas dans sa bouche.

\pagerange{3}{5}
\section*{Le groupe $S_{E}$ et les classes d'idèles}

Les trois dernières cartes changent de terrain. Il ne s'agit plus d'une place,
mais d'un corps de nombres $E$ et du groupe algébrique par lequel passent ses
caractères de Hecke algébriques.

\subsection*{La norme}

La troisième carte définit une flèche
\[
N : S_{E} \longrightarrow S_{\mathbb{Q}} \longrightarrow \mathbb{G}_{m},
\qquad
(x, (y_{\alpha})) \longmapsto x \cdot \prod_{\alpha \neq \infty}
\lVert y_{\alpha}\rVert \cdot \prod_{\alpha = \infty}
\mathrm{sign}(y_{\alpha}) .
\]
La dissymétrie de la formule est son contenu : aux places finies on prend la
valeur absolue, à la place infinie on ne prend que le signe. Elle est ensuite
composée dans la chaîne
\[
C(E) \longrightarrow S_{E}(\mathbb{A}) \longrightarrow S_{E}(\mathbb{R})
\longrightarrow \mathbb{R}^{*},
\]
où $C(E)$ est le groupe des classes d'idèles, et la carte note sous cette
composée, souligné deux fois, le symbole $\lVert x\rVert$ : c'est donc la norme
adélique qu'on retrouve au bout.

En bas de la carte, isolée sous un trait, une seule remarque :
\[
(-1 \neq \lVert -1 \rVert).
\]
Elle n'est pas anodine. Elle dit que la flèche construite n'est pas la valeur
absolue : elle garde le signe là où la valeur absolue le perd, et c'est
précisément cet écart qui laisse place au facteur $\gamma_{\mathbb{R}}$ décalé
d'une unité rencontré plus haut.\footnote{Le rapprochement entre cette remarque
et la règle de la place réelle est le nôtre. Les cartes sont contiguës dans le
dossier et parlent du même objet, mais aucune ne renvoie explicitement à
l'autre.}

\subsection*{Le diagramme}

La cinquième carte rassemble le tout en un diagramme, sous un $E/\mathbb{Q}$
encadré :

\begin{tikzcd}[column sep=small, row sep=small, nodes={font=\scriptsize}]
& & & \mathbb{G}_{m}\\
E^{*} \arrow[r] & \left(\prod_{E/\mathbb{Q}} \mathbb{G}_{m}\right)\big/U \times I(E)/I^{\circ}(E) \arrow[rr] & & S_{E} \arrow[u, "N"]\\
E^{*} \arrow[r] \arrow[u, no head] & I(E) \times I(E)/I^{\circ}(E) \arrow[rr] & & S_{E}(\mathbb{A})\\
E^{*} \arrow[rr] \arrow[u, no head] & & I(E) \arrow[r] & C(E) \arrow[u]
\end{tikzcd}

Les trois lignes se lisent de bas en haut comme un raffinement progressif. En
bas, la construction ordinaire du groupe des classes d'idèles : $E^{*}$ se
plonge dans le groupe des idèles $I(E)$, et le quotient est $C(E)$. Au milieu,
on adjoint le groupe des classes de rayon $I(E)/I^{\circ}(E)$, ce qui donne les
points adéliques de $S_{E}$. En haut, la restriction des scalaires
$\prod_{E/\mathbb{Q}}\mathbb{G}_{m}$ remplace les idèles et l'on atteint $S_{E}$
lui-même, sur lequel la norme $N$ de la carte précédente est
définie.\footnote{Les trois traits verticaux de la colonne de gauche sont, sur
la page, des doubles barres d'égalité : c'est le même $E^{*}$ aux trois lignes.
On les rend ici par des liaisons sans tête, la syntaxe des diagrammes de cette
édition n'ayant pas de barre d'égalité. Le $S_{E}(\mathbb{Q})$ que la carte
écrit au-dessus de $S_{E}(\mathbb{A})$ est barré et n'est pas reporté ; deux
flèches obliques numérotées 1 et 2 courent entre la deuxième et la troisième
ligne sans origine assurée, et ne le sont pas non plus.}

Sous un trait, en bas de la carte, une flèche $S_{E} \to \mathbb{G}$ est écrite
puis barrée, et remplacée par $S_{E} \to \mathbb{G}_{m}$ --- la norme, donc,
reprise une dernière fois. La suite sort de la carte.\footnote{On y distingue
un $h$ et une parenthèse ouvrante suivie d'un $E$ ; ce pourrait être un nombre
de classes, mais la restitution serait la nôtre et on ne la propose pas.}

\pagerange{1}{5}
\section*{Ce que les cinq cartes font ensemble}

Prises séparément, ce sont cinq notes sans énoncé. Prises ensemble, elles
dessinent un aller-retour.

D'un côté, la géométrie : une variété, sa cohomologie, ses nombres de Hodge,
et la recette qui en tire le facteur à l'infini. De l'autre, l'arithmétique :
un corps de nombres, ses idèles, et le groupe $S_{E}$ qui gouverne ses
caractères. Entre les deux, le groupe de Weil, qui a la particularité d'avoir
les deux natures --- ses représentations sont du côté arithmétique, et ce
qu'elles classent est du côté de Hodge.

La dernière carte du dossier, la quatrième, écrit d'ailleurs le seul énoncé de
compatibilité de tout l'ensemble :
\[
L_{v}(M \otimes T, s) = L_{v}(M)(s \pm 1),
\]
où $T$ est l'objet de Tate. Elle dit que tordre un motif par $T$ revient à
translater sa fonction $L$ d'une unité --- et le $\pm$ n'est pas une hésitation
sur le fait, mais sur la convention de signe adoptée pour le
twist.\footnote{Cette lecture du $\pm$ est la nôtre. La carte l'écrit sans
commentaire, et l'on pourrait aussi bien y voir une hésitation sur la
normalisation de $T$ lui-même. Aucune autre carte du dossier ne tranche.}

\end{document}
